\documentclass{ximera}
\title{Exponential Growth and Decay}
\begin{document}
\begin{abstract}
Quantities that grow or decay in proportion to themselves, the exponential solutions $y(t)=y(0)e^{kt}$, radioactive decay and Newton's law of cooling. Textbook §3.8.
\end{abstract}
\maketitle

We've also done some exponential growth and decay problems already. We saw this in the population growth case.

A lot of times, things grow and decay proportional to themselves. So in essence, in a population, you grow based on how many people are present, and that growth in measured by some proportion of the current number. As an example: If we take a census in 2015 and then take a census in 2020 then we might find that the population grew by 20\%. What this means is that if $P$ is the population in 2015, then $1.2 P$ is the population in 2020. So the population is measured in proportions.

This rate of change can be made instantaneous by taking the derivative, giving us:
\[
\frac{d y}{d t}=k y
\]
where $k$ is the proportion.

We've only seen one function where $y$ stays itself after taking the derivative! The exponential!!! That means, usually with growth and decay, we are dealing with the exponential function.

In fact, it is known that the only solutions to the equation $\frac{d y}{d t} k y$ are the exponential functions:
\[
y(t)=y(0) e^{k t}
\]

Let's look at some examples where we use these exponential functions.

\begin{example}\label{example:8-3}
Let's first look at the nucleur armageddon. In other words: radioactive decay. This example is different from the population growth one because we are actually decreasing the volume of a radioactive substance rather than increasing. In essence, if we have some particle then it's mass decays exponentially by the following formula:
\[
m(t)=m_{0} e^{k t}
\]
where $m_{0}$ is the initial value, $k$ is some constant for the particle, and $t$ is the variable. In order to calculate this, physicists use something called the half-life of a particle. The half-life of a particle is the time needed for half of the particle to decay. We use all these facts in the following example.

So let's suppose we have $1 kg$ of uranium-235. It takes 704 million years for 1 gram of uranium-235 to decay to 1/2 a gram. In other words, the half-life of uranium-235 is 704 million years $=704000000$. We already know that we are starting off with 1000 grams of uranium-235 so that means $m_{0}=1000$. But how do we find $k$?

Using the half-life! We already know that $m(704000000)$ should be half the initial mass since 704000000 is the half-life. So So
% work: 500 = 1000 e^{704000000 k}  =>  e^{704000000 k} = 1/2  =>  k = ln(1/2)/704000000 = -ln(2)/704000000
\[
m(704000000)=1000e^{704000000k}=500
\quad\Longrightarrow\quad
k=\answer{-\frac{\ln 2}{704000000}}
\]
\[
m(t)=\answer{1000e^{-\frac{\ln 2}{704000000}t}}
\]

How much of our uranium will be left over in 5 billion years when the sun dies out?
% work: m(5000000000) = 1000 e^{-ln(2) 5000000000/704000000} = 1000 * 2^{-5000/704} ≈ 7.28 grams
\[
m(5000000000)=\answer{1000\cdot 2^{-5000000000/704000000}}\text{ grams}
\]
\end{example}

\begin{example}\label{example:8-4}
Our next example uses Newton's law of cooling. We're going to be looking at the temperature of an object over some time, so we denote this function by $T(t)$. Newton's law of cooling says that
\[
\frac{d T}{d t}=k\left(T-T_{s}\right)
\]
where $T_{s}$ is the temperature of the surroundings. To calculate things, we need to fix things up. So we let $y(t)=T(t)-T_{s}$. Then $y^{\prime}(t)=T^{\prime}(t)$ and so our equation above becomes:
\[
y^{\prime}(t)=k y
\]
which matches our exponential formula from above.

So suppose that we just made some hot boiling tea (or coffee) and it's just chilaxing at room temperature ($20^{\circ}$). We know that in 30 minutes, the tea will be $55^{\circ}$. Considering that the perfect drinking temperature for tea is $65^{\circ}$, when should you drink your tea?

This is a ton of information! So let's break it down slow. First, let's figure out our formula:
% work: T_s = 20 (room temperature)
\[
\frac{d T}{d t}=k\left(T-T_{s}\right)=\answer{k(T-20)}
\]
% work: boiling tea: T(0) = 100, y(0) = 100 - 20 = 80; T(30) = 55, y(30) = 55 - 20 = 35
We also know that $T(0)=\answer{100}$ and so $y(0)=T(0)-T_{s}=\answer{80}$.
Also, we know that $T(30)=\answer{55}$ and so $y(30)=T(30)-T_{s}=\answer{35}$.

Then $y(t)=y(0) e^{k t}$ (from our equation earlier!). We can now calculate $k$:
% work: y(30) = y(0) e^{30k}  =>  35 = 80 e^{30k}  =>  k = (1/30) ln(35/80) ≈ -0.02756
\[
y(30)=y(0)e^{30k}
\quad\Longrightarrow\quad
k=\answer{\frac{1}{30}\ln\left(\frac{35}{80}\right)}
\]

So now we want to figure out when $T(t)=\answer{65}$ in other words when $y(t)=\answer{45}$. Therefore:
% work: 45 = 80 e^{kt}  =>  t = ln(45/80)/k = 30 ln(45/80)/ln(35/80) ≈ 20.88 minutes
\[
y(t)=y(0)e^{kt}
\quad\Longrightarrow\quad
t=\answer{\frac{30\ln\left(\frac{45}{80}\right)}{\ln\left(\frac{35}{80}\right)}}
\]
In other words, we must wait about 20 minutes after pouring tea before it is at perfect temperature.
\end{example}

\end{document}
